\emph{Preliminaries.} The adjoint, the scalar part and their properties ($\langle x^{\dagger}\rangle_{0} = \langle x\rangle_{0}$, $\langle xy\rangle_{0} = \langle yx\rangle_{0}$, $\langle x^{\dagger}x\rangle_{0}$ the sum of the squares of the coefficients of $x$, the scalar part of a product of two blades zero unless they are equal up to sign) are those of Theorem \claimref{synthesis.chsh-bound-for-two-subsystems-of-a-clifford-algebra-by-a-rational-certificate}; the blades are a basis by Theorem \claimref{synthesis.signed-directions-and-the-integer-clifford-algebra}. By clause (vii) of the former and Theorem \claimref{tower.commuting-copies}, $X_{1}, Z_{1}, X_{2}, Z_{2}$ are self-adjoint blades of square $1$, $X_{1}Z_{1} = -Z_{1}X_{1}$, $X_{2}Z_{2} = -Z_{2}X_{2}$, and each of $X_{1}, Z_{1}$ commutes with each of $X_{2}, Z_{2}$. Hence a direction has square $x^{2} + z^{2} = 1$ and is self-adjoint, and directions of different parties commute. $U = X_{1}X_{2}$ and $V = Z_{1}Z_{2}$ are products of commuting self-adjoint elements of square $1$, so they are self-adjoint of square $1$; $UV = X_{1}X_{2}Z_{1}Z_{2} = X_{1}Z_{1}X_{2}Z_{2} = VU$ (two anticommutations), and $X_{1}Z_{1} = e_{1}f_{1}e_{1} = -f_{1}$, $X_{2}Z_{2} = e_{2}p_{1}f_{2}e_{2} = -f_{2}p_{1}$, so $UV = f_{1}f_{2}p_{1} = f_{1}f_{2}f_{1}e_{1} = f_{2}e_{1} = -K$ with $K = e_{1}f_{2}$, a blade with $K^{\dagger} = f_{2}^{\dagger}e_{1}^{\dagger} = -f_{2}e_{1} = K$ and $K^{2} = 1$. The nine products spanning $L$ are, up to sign, the nine distinct blades $1, e_{1}, p_{1}, e_{2}p_{1}, p_{2}, e_{1}e_{2}p_{1}, e_{1}p_{2}, p_{1}e_{2}p_{1} = e_{2}, p_{1}p_{2}$, none of them $\pm e_{1}f_{2}$; so $K \notin L$, while $K = -(X_{1}Z_{1})(X_{2}Z_{2})$ with $X_{1}Z_{1} = -f_{1}$ and $X_{2}Z_{2} = -f_{2}p_{1}$ anti-self-adjoint (verified by the check), and $\langle Ky\rangle_{0} = 0$ for $y \in L$.

\emph{(i).} $y = \tfrac14(1 + a)(1 + b)$ is a product of two commuting self-adjoint idempotents, so $y^{\dagger} = y = y^{2} = yy^{\dagger}$, and (P) gives $\lambda(y) \ge 0$.

\emph{(ii).} If $a$, $b$ are directions so are $-a$, $-b$. As elements, $\tfrac14(1 + a)(1 + b) + \tfrac14(1 - a)(1 - b) = \tfrac12(1 + ab)$ and $\tfrac14(1 + a)(1 - b) + \tfrac14(1 - a)(1 + b) = \tfrac12(1 - ab)$. Under (L) the left sides have non-negative values, so $1 + \lambda(ab) \ge 0$ and $1 - \lambda(ab) \ge 0$: $|\lambda(ab)| \le 1$. The map $\beta(a, b) = \lambda(ab)$, extended linearly to the plane spanned by $X_{1}, Z_{1}$ and the plane spanned by $X_{2}, Z_{2}$, is a bilinear form on these two planes with values in $k \subseteq \R$, fixed by its four values on the generators; write $|a|^{2} = x^{2} + z^{2}$ for $a = xX_{1} + zZ_{1}$, which is $a^{2}$, and likewise for $b$. So $|\beta(a, b)| \le 1$ for all $a$, $b$ of norm one with coordinates in $k$. The points of norm one with coordinates in $\Q$, hence in $k$, are dense in those with real coordinates (the points $((1 - s^{2})/(1 + s^{2}), 2s/(1 + s^{2}))$, $s \in \Q$), and $\beta$, extended to real coordinates, is continuous, so $|\beta(a, b)| \le 1$ at norm one and $|\beta(a, b)| \le |a|\,|b|$ for all elements of the two real planes. For directions $a_{1}, a_{2}, b_{1}, b_{2}$,
\[
\lambda(C) = \beta(a_{1}, b_{1} + b_{2}) + \beta(a_{2}, b_{1} - b_{2}), \qquad |\lambda(C)| \le |b_{1} + b_{2}| + |b_{1} - b_{2}| ,
\]
and $|b_{1} + b_{2}|^{2} + |b_{1} - b_{2}|^{2} = 2|b_{1}|^{2} + 2|b_{2}|^{2} = 4$, so $\lambda(C)^{2} \le 2 \cdot 4 = 8$ by $(s + d)^{2} \le 2(s^{2} + d^{2})$. For $k = \Q$ the value $\lambda(C)$ is rational and $8$ is not the square of a rational, so the inequality is strict. Only the weights at pairs of directions were used.

\emph{(iii).} $C_{0} = X_{1}X_{2} + X_{1}Z_{2} + Z_{1}X_{2} - Z_{1}Z_{2}$ is a sum of four distinct blades up to sign, none scalar, so $\lambda_{\mathrm{box}}(1) = 1$, $\lambda_{\mathrm{box}}(C_{0}) = \langle C_{0}^{2}\rangle_{0} = 4$, the values on the four products are $1, 1, 1, -1$, and the values on $X_{1}, Z_{1}, X_{2}, Z_{2}$ are $0$; this is the functional of Theorem \claimref{synthesis.chsh-bound-needs-positivity-across-settings-and-anticommutation-holds-only-on-the-state} with the settings named in another order. For $b = \tfrac35X_{2} + \tfrac45Z_{2}$: $\lambda_{\mathrm{box}}(X_{1}b) = \tfrac35 + \tfrac45 = \tfrac75$ and $\lambda_{\mathrm{box}}\bigl(\tfrac14(1 + X_{1})(1 - b)\bigr) = \tfrac14(1 - \tfrac75) = -\tfrac{1}{10}$.

\emph{(iv).} $M_{c} = 1 - U - V + cK$ is self-adjoint and even, with scalar part $1$, so $\lambda_{c}(1) = 1$, $\lambda_{c}(x^{\dagger}) = \langle M_{c}x^{\dagger}\rangle_{0} = \langle xM_{c}\rangle_{0} = \lambda_{c}(x)$, and $\lambda_{c}$ vanishes on odd $x$ (the product $M_{c}x$ is then odd). On $L$ the term $cK$ contributes nothing, by the preliminaries; the other terms give $\lambda_{c}(X_{1}X_{2}) = \lambda_{c}(Z_{1}Z_{2}) = -1$ and $0$ on the other seven products, whence $\lambda_{c}(ab) = -(xx' + zz')$, $\lambda_{c}(a) = \lambda_{c}(b) = 0$ and the weight $\tfrac14(1 - xx' - zz') \ge 0$ (Cauchy's inequality for two unit vectors): (L) holds for every $c$. $\lambda_{c}(K) = c\langle K^{2}\rangle_{0} = c$. Since $U$, $V$ commute, have square $1$ and $UV = -K$, the four elements $\tfrac14(1 + uU)(1 + vV)$ are self-adjoint idempotents and $M_{c}\,\tfrac14(1 + uU)(1 + vV) = (1 - u - v - c\,uv)\,\tfrac14(1 + uU)(1 + vV)$; each has scalar part $\tfrac14$ ($U$, $V$, $UV$ are non-scalar blades), which gives the table. If $c \ne -1$ one of $\mp\tfrac14(1 + c)$ is negative on an element $yy^{\dagger}$, so (P) fails. If $c = -1$, $M_{-1} = (1 - U)(1 - V) = 4P_{s}$ and $\lambda_{-1}(xx^{\dagger}) = 4\langle P_{s}xx^{\dagger}P_{s}\rangle_{0} = 4\langle(P_{s}x)(P_{s}x)^{\dagger}\rangle_{0} \ge 0$, so (P) holds. The identity $P_{s}xP_{s} = 4\langle P_{s}x\rangle_{0}P_{s}$ is linear in $x$ and is verified on the sixteen blades by the check; so $P_{s}$ is a projector with corner $kP_{s}$, an atom by the Definitions of Theorem \claimref{synthesis.born-weights-under-positivity-a-state-sharp-on-an-atom-is-its-trace-functional}, $\lambda_{-1} = \omega_{P_{s}}$, and uniqueness of the state sharp on $P_{s}$ is clause (iii) of that theorem. $\lambda_{c}(P_{s}) = \tfrac14\langle M_{c}(1 - U - V - K)\rangle_{0} = \tfrac14(3 - c)$. The check also evaluates every statement of this clause for five values of $c$; each is affine in $c$.

\emph{(v).} The average over conjugation by one mode parity is the descent of Theorem \claimref{singularity.descent-parity-average} for that mode, and the composite of the mode-wise averages is the map $E_{N}$ of Proposition \claimref{singularity.p-reduced-metric-level-reading-many}; the same theorem with $\Gamma$ in place of the mode parity is the total-parity average $E$. Conjugation by $p_{j}$ multiplies a blade by $\pm 1$, so $E_{N}$ is the projection onto the span of the blades fixed by both conjugations, which are $1, p_{1}, p_{2}, p_{1}p_{2}$; this span is also the span of the four $\Pi_{n}$, orthogonal projectors with sum $1$. That $\lambda \circ E_{N}$ satisfies (P) exactly when the four $\lambda(\Pi_{n})$ are non-negative, and that the total-parity average adds no positivity, is the rider and the statement of clause (iv) of Theorem \claimref{synthesis.positivity-of-values-is-additivity-of-non-negative-weights-on-effects}; the check verifies the formula $E_{N}(x) = \sum_{n}\Pi_{n}x\Pi_{n} = 4\sum_{n}\langle\Pi_{n}x\rangle_{0}\Pi_{n}$ on all blades for the two-mode instances. For $M_{c}$ and $1 + C_{0}$ the fixed part is $1 - V$ ($V = p_{1}p_{2}$; $U$, $K$ and the other blades of $C_{0}$ are not fixed), with weights $0, \tfrac12, \tfrac12, 0$, and $1 - V = \tfrac12(1 - V)^{2}$. For the total parity, conjugation by $\Gamma$ fixes even blades and negates odd ones, so $M_{c}$ is fixed and the even part of $1 + C_{0}$ is $1 + U - V$; with $y = \tfrac14(1 - U)(1 + V) = yy^{\dagger}$, $(1 + U - V)y = -y$ and the value is $-\langle y\rangle_{0} = -\tfrac14$.
